% =============================================================================
% Estilo de la memoria: sobrio y academico (sin color, tipografia Times).
% Aqui solo hay DISEÑO. El contenido esta en secciones/, las figuras en figuras/.
%
% Lo que se usa al escribir (ver ejemplos en secciones/):
%   \pendiente{...}        texto gris en cursiva: lo que falta por escribir (borrar al redactar)
%   \falta{...}            lo mismo, dentro de una linea (tablas, trazas)
%   \paragraph{Nombre.} ...                  parrafo de un algoritmo (titulo en negrita en linea)
%   algorithmic (paquete algpseudocode)       pseudocodigo, con palabras clave en castellano
%   \begin{traza}{titulo} ... \end{traza}    traza de ejecucion (texto monoespaciado)
%   \tier{n} \est{ok|aprox|pend}             casillas de la tabla de tareas (3_tareas.tex)
%   \puntuacion                              puntos calculados con las casillas "ok"
% =============================================================================

% ---------------------------------------------------------------- Fuentes
\usepackage[T1]{fontenc}
\usepackage[spanish,es-tabla]{babel}
\usepackage{amsmath}
\usepackage{newtxtext,newtxmath}        % Times para texto y matematicas
\usepackage[scaled=0.9]{zi4}            % monoespaciada (Inconsolata) para trazas
\usepackage{microtype}                  % mejor reparto de espacios: ahorra lineas

% ---------------------------------------------------------------- Pagina
% Enunciado: max. 5 paginas, fuente 12 pt, margenes <= 2,54 cm, interlineado simple.
\usepackage[a4paper,margin=2.2cm]{geometry}
\setlength{\parindent}{1em}
\setlength{\parskip}{0pt}
\pagestyle{plain}                       % solo el numero de pagina, abajo

% ---------------------------------------------------------------- Paquetes
\usepackage{booktabs,tabularx,array,multirow}
\usepackage{xcolor}
\usepackage{graphicx}
\usepackage{float}                      % [H]: la tabla o figura va justo donde se escribe
\usepackage{enumitem}
\setlist{nosep,leftmargin=1.5em}
\usepackage{tikz}
\usetikzlibrary{babel,arrows.meta,positioning,calc}
\usepackage{pgfplots}
\pgfplotsset{compat=1.18}
\usepackage[font=small,labelfont=bf,labelsep=period]{caption}
\usepackage{algorithm,algpseudocode}    % pseudocodigo (Algoritmo 1: LLL)
\usepackage[hidelinks]{hyperref}

\definecolor{gris}{gray}{0.45}          % unico color: marcas de "pendiente"

% ---------------------------------------------------------------- Titulos de seccion
\usepackage{titlesec}
\titleformat{\section}{\normalfont\large\bfseries}{\thesection}{0.6em}{}
\titleformat{\subsection}{\normalfont\normalsize\bfseries}{\thesubsection}{0.6em}{}
\titlespacing*{\section}{0pt}{10pt plus 2pt minus 2pt}{4pt}
\titlespacing*{\subsection}{0pt}{6pt plus 2pt minus 1pt}{2pt}

% ---------------------------------------------------------------- Titulo
% \cabecera{titulo}{autores}{asignatura}: bloque de titulo centrado.
\newcommand{\cabecera}[3]{%
  \begin{center}
    {\LARGE\bfseries #1\par}\vspace{8pt}
    {\large #2\par}\vspace{3pt}
    {\small\itshape #3\par}
  \end{center}\vspace{-2pt}}

% ---------------------------------------------------------------- Pendientes
% Antes de entregar no debe quedar ninguno: buscar "\pendiente" y "\falta" en secciones/.
\newcommand{\pendiente}[1]{\par{\small\itshape\color{gris}[Pendiente: #1]\par}}
\newcommand{\falta}[1]{\leavevmode{\normalfont\small\itshape\color{gris}[#1]}}

% ---------------------------------------------------------------- Parrafos de algoritmo
% \paragraph{Babai.} texto...  -> titulo en negrita al principio de la linea, sin numero.
\titleformat{\paragraph}[runin]{\normalfont\normalsize\bfseries}{}{0pt}{}
\titlespacing*{\paragraph}{0pt}{3pt plus 1pt}{0.5em}

% ---------------------------------------------------------------- Pseudocodigo en castellano
\floatname{algorithm}{Algoritmo}
\captionsetup[algorithm]{singlelinecheck=false}   % titulo del algoritmo alineado a la izquierda
\algrenewcommand\algorithmicrequire{\textbf{Entrada:}}
\algrenewcommand\algorithmicensure{\textbf{Salida:}}
\algrenewcommand\algorithmicwhile{\textbf{mientras}}
\algrenewcommand\algorithmicfor{\textbf{para}}
\algrenewcommand\algorithmicdo{\textbf{hacer}}
\algrenewcommand\algorithmicif{\textbf{si}}
\algrenewcommand\algorithmicthen{\textbf{entonces}}
\algrenewcommand\algorithmicelse{\textbf{si no}}
\algrenewcommand\algorithmicend{\textbf{fin}}
\algrenewcommand\algorithmicreturn{\textbf{devolver}}
\algrenewcommand\algorithmiccomment[1]{\hfill{\itshape\footnotesize #1}}

% ---------------------------------------------------------------- Traza de ejecucion
% Resultados intermedios de una ejecucion real (lo pide el enunciado), entre dos filetes.
%   \begin{traza}{Tier0T01\_Latt (LLL + Babai)}
%   base reducida: ...\\
%   \end{traza}
\newenvironment{traza}[1]
  {\par\smallskip\noindent{\small\textbf{Traza.} #1}\par\nopagebreak
   \noindent\rule{\linewidth}{0.4pt}\par\nopagebreak\vspace{1pt}
   \noindent\begin{minipage}{\linewidth}\ttfamily\small}
  {\end{minipage}\par\nopagebreak\vspace{1pt}\noindent\rule{\linewidth}{0.4pt}\par\smallskip}

% ---------------------------------------------------------------- Tabla de tareas
% \est{estado}: ok = resuelta y comprobada, aprox = resuelta sin certificar el optimo,
% pend = pendiente. Las "ok" se cuentan por tier (\tier{n} al principio de cada fila)
% y \puntuacion calcula los puntos con los maximos del enunciado.
\newcounter{oktierA}\newcounter{oktierB}\newcounter{oktierC}
\newcommand{\tieractual}{0}
\newcommand{\tier}[1]{\gdef\tieractual{#1}}
\newcommand{\est}[1]{%
  \ifnum\pdfstrcmp{#1}{ok}=0
    \ifnum\tieractual=1 \stepcounter{oktierA}\fi
    \ifnum\tieractual=2 \stepcounter{oktierB}\fi
    \ifnum\tieractual=3 \stepcounter{oktierC}\fi
    $\checkmark$%
  \else\ifnum\pdfstrcmp{#1}{aprox}=0 $\approx$%
  \else --\fi\fi}
% Tier 1: 0,3 pts (max 1,4); Tier 2: 0,5 (max 1,8); Tier 3: 1 (sin max); total max 4.
\newcommand{\puntuacion}{%
  \pgfmathsetmacro{\pA}{min(0.3*\value{oktierA},1.4)}%
  \pgfmathsetmacro{\pB}{min(0.5*\value{oktierB},1.8)}%
  \pgfmathsetmacro{\pC}{1.0*\value{oktierC}}%
  \pgfmathsetmacro{\pT}{min(\pA+\pB+\pC,4)}%
  \def\num##1{\pgfmathprintnumber[fixed,fixed zerofill,precision=1,use comma]{##1}}%
  Tier~1: \num{\pA}; Tier~2: \num{\pB}; Tier~3: \num{\pC}; \textbf{total: \num{\pT}\,/\,4}}

% ---------------------------------------------------------------- Tablas
\newcolumntype{Y}{>{\raggedright\arraybackslash}X}

% ---------------------------------------------------------------- Aviso de longitud
% Si la memoria pasa de 5 paginas, aparece un aviso al final (y en el log).
\AtEndDocument{%
  \ifnum\value{page}>5
    \par\medskip\noindent{\small\itshape\color{gris}%
      [Aviso: la memoria tiene \thepage\ páginas; el máximo son 5.]}%
    \PackageWarningNoLine{memoria}{La memoria pasa de 5 paginas}%
  \fi}
